\begin{table}[htbp]\centering
\begin{threeparttable}
\caption{The complex rule applied as a prediction within each round of the LSU-removal test.}
\label{tab:foldwise}
\footnotesize\setlength{\tabcolsep}{2pt}
\begin{tabular}{lrcrlrr}
\toprule
Withheld species & Sequences & Complex & \multicolumn{2}{c}{Inside a remaining congener} & \multicolumn{2}{c}{False known}\\
 & & (all cultures) & $n$ & epithet ($n$) & no rule & ITS rule\\
\midrule
\emph{Ambispora leptoticha} & 10 & -- & 0 & -- & 10 & 0\\
\emph{Diversispora aestuarii} & 20 & C2 & 20 & \emph{varaderana} (20) & 20 & 20\\
\emph{Diversispora arenaria} & 4 & C3 & 4 & \emph{jakucsiae} (4) & 4 & 4\\
\emph{Diversispora jakucsiae} & 4 & C3 & 4 & \emph{arenaria} (4) & 4 & 4\\
\emph{Diversispora varaderana} & 31 & C2 & 31 & \emph{aestuarii} (31) & 3 & 3\\
\emph{Funneliformis geosporum} & 8 & -- & 0 & -- & 8 & 0\\
\emph{Glomus macrocarpum} & 7 & C1 & 7 & \emph{rugosae} (7); \emph{chinense} (1) & 4 & 4\\
\emph{Rhizoglomus irregulare} & 56 & C4 & 13 & \emph{vesiculiferum} (13) & 19 & 18\\
\emph{Septoglomus africanum} & 6 & -- & 0 & -- & 6 & 0\\
\emph{Rhizoglomus vesiculiferum} & 23 & C4 & 2 & \emph{irregulare} (2) & 0 & 0\\
\emph{Glomus chinense}$^{\dagger}$ & 35 & C1 & 6 & \emph{rugosae} (6) & 28 & 5\\
\emph{Glomus rugosae}$^{\dagger}$ & 75 & C1 & 62 & \emph{macrocarpum} (62); \emph{chinense} (9) & 75 & 62\\
24 other test species & 504 & -- & 0 & -- & 0 & 0\\
\bottomrule
\end{tabular}
\begin{tablenotes}\footnotesize
\item Main-type sequences of each withheld species; $n$, number of sequences. \emph{Complex}: the species' complex when the rule is computed from all
cultures (Table~\ref{tab:complexes}). \emph{Inside}: with the withheld cultures removed and the within-species limit
recomputed without them, sequences that lie within the limit of a remaining species on all three markers; only these
can be assigned to a known species if the rule holds. False-known assignments with the species removed from the LSU reference
(Table~\ref{tab:m4folds}). $^{\dagger}$\,Development species.
\end{tablenotes}
\end{threeparttable}
\end{table}
